\section{Executive Summary}

This assessment evaluates the feasibility and implications of migrating the organization's core transaction processing system from a monolithic architecture to a microservices-based event-driven architecture implementing the Command Query Responsibility Segregation (CQRS) pattern. The decision requires balancing immediate operational complexity against long-term scalability requirements as transaction volumes approach the current system's design limits.

The current monolithic system processes approximately 45,000 transactions per second at peak load, with p99 latency of 280ms and p99.9 latency of 1,200ms. Capacity modeling indicates the system will reach hard scalability limits at 65,000 TPS, projected to occur within 18 months based on current growth trajectories. Vertical scaling options have been exhausted; the existing database cluster operates at 78\% of maximum IOPS across all availability zones.

The proposed microservices architecture introduces 12 bounded contexts, each deployable independently, with event sourcing for state changes and CQRS for read-write separation. Benchmark data from pilot implementations demonstrate horizontal scalability to 200,000+ TPS with maintained p99 latency below 150ms. However, this performance comes at substantial cost: operational complexity increases by approximately 4x as measured by deployment artifacts, monitoring surface area, and failure mode enumeration.

The migration path spans 24 months with an estimated total cost of \$8.2M including engineering time, infrastructure provisioning, and opportunity cost of delayed features. The first 6 months focus on platform preparation: service mesh deployment, observability infrastructure, and chaos engineering capabilities. Months 7-18 involve incremental strangler pattern migration of individual bounded contexts. The final 6 months address data consistency verification and cutover procedures.

Critical risks include: (1) underestimating inter-service transaction complexity, with 23\% of current operations spanning multiple proposed service boundaries; (2) operational readiness gaps, as current SRE team size supports 4 production services versus the proposed 12; (3) data migration integrity, particularly for the 400TB transactional history that must remain queryable during migration.

The assessment recommends conditional approval contingent on three prerequisites: (a) completion of organizational readiness program to expand SRE capacity; (b) successful pilot of saga pattern for the three highest-complexity cross-service transactions; (c) vendor selection for managed Kubernetes and service mesh to reduce operational burden. With these prerequisites met, the architecture change addresses the scalability constraint while maintaining acceptable risk levels.
